\chapter{Chapter 2 Ansys Fluent Benchmark Support}
\label{app:A_backstep_benchmark}

This appendix supports the turbulence-modeling discussion in Chapter~\ref{ch:turbulence_modeling}. It contains the common Ansys Fluent convergence-monitor definitions for the Chapter~\ref{ch:turbulence_modeling} benchmark runs and the backward-facing-step material that is useful as separated-flow model-sensitivity evidence. The backward-facing-step case is secondary to the main Chapter~\ref{ch:turbulence_modeling} argument based on the Driver--Seegmiller/NASA literature comparison and is not used as a FEniCS solver-verification benchmark.

Unless stated otherwise, all grid-sensitivity relative differences in this appendix are reported as positive percentage magnitudes. For a mesh-dependent metric $\phi$, the coarse--medium difference is computed as $|\phi_{\mathrm{coarse}}-\phi_{\mathrm{medium}}|/|\phi_{\mathrm{coarse}}|\times100\%$, while the medium--fine difference is computed as $|\phi_{\mathrm{medium}}-\phi_{\mathrm{fine}}|/|\phi_{\mathrm{fine}}|\times100\%$.

\section{Fluent Benchmark Convergence Monitors}
\label{sec:app_A_ch2_fluent_convergence}

For the Fluent benchmark runs in Section~\ref{sec:ansys_benchmarks}, convergence was assessed using Fluent report monitors. Three physical monitors were tracked in addition to the solver residuals: the inlet-to-outlet static pressure drop, the domain-integrated viscous dissipation, and the relative mass imbalance. The pressure-drop monitor was defined from area-weighted static pressures,
\begin{equation}
	\Delta p = \bar{p}_{\mathrm{in}} - \bar{p}_{\mathrm{out}},
	\label{eq:ch2_fluent_pressure_drop_monitor}
\end{equation}
while the viscous-dissipation monitor used a Fluent custom field function,
\begin{equation}
	\Phi_D = \int_{\Omega} \mu_{\mathrm{eff}}\dot{\gamma}^{2} \, d\Omega ,
	\label{eq:ch2_fluent_viscous_dissipation_monitor}
\end{equation}
where the custom field function \texttt{viscous-dissipation-density} was defined as $\mu_{\mathrm{eff}}\dot{\gamma}^{2}$ using Fluent's \texttt{effective-viscosity} and \texttt{strain-rate-magnitude} variables. Here $\mu_{\mathrm{eff}}=\mu+\mu_t$ for turbulent runs and $\dot{\gamma}$ is Fluent's scalar strain-rate magnitude. For a monitored scalar $q$, the final-window change over $N$ iterations was evaluated as
\begin{equation}
	\delta_N(q) =
	\frac{\left|q^{(k)} - q^{(k-N)}\right|}
	{\left|q^{(k)}\right|} .
	\label{eq:ch2_fluent_monitor_window_change}
\end{equation}
The relative mass imbalance was computed from the signed inlet and outlet mass-flow-rate reports,
\begin{equation}
	\epsilon_m =
	\frac{\left|\dot{m}_{\mathrm{in}} + \dot{m}_{\mathrm{out}}\right|}
	{\left|\dot{m}_{\mathrm{out}}\right|},
	\label{eq:ch2_fluent_mass_imbalance}
\end{equation}
where $\dot{m}_{\mathrm{out}}<0$ under Fluent's boundary-normal sign convention for these cases.

\begin{table}[!ht]
	\centering
	\scriptsize
	\caption{Convergence metrics used for the Ansys Fluent benchmark runs in Section~\ref{sec:ansys_benchmarks}. The residual criterion refers to all Fluent-scaled residuals, including continuity, momentum, and turbulence-transport equations.}
	\label{tab:ch2_ansys_convergence_metrics}
	\renewcommand{\arraystretch}{1.2}
	\begin{tabularx}{\textwidth}{X c c c c c}
		\hline
		\textbf{Benchmark} &
		\textbf{Window} &
		\shortstack{\textbf{Pressure-drop}\\\textbf{change}\\$\delta_N(\Delta p)$} &
		\shortstack{\textbf{Dissipation}\\\textbf{change}\\$\delta_N(\Phi_D)$} &
		\shortstack{\textbf{Mass}\\\textbf{imbalance}\\$\epsilon_m$} &
		\shortstack{\textbf{Scaled}\\\textbf{residuals}} \\
		\hline
		VMFL030 $90^\circ$ bend, 3D & $500$ iters & $<10^{-3}$ & $<10^{-3}$ & $<10^{-3}$ & $10^{-6}$ \\
		VMFL048 U-bend, 3D & $500$ iters & $<10^{-3}$ & $<10^{-3}$ & $<10^{-3}$ & $10^{-6}$ \\
		Backward-facing step, 2D & $250$ iters & $<10^{-3}$ & $<10^{-3}$ & $<10^{-3}$ & $10^{-6}$ \\
		\hline
	\end{tabularx}
\end{table}

\section{Ansys SA--\texorpdfstring{SST $k-\omega$}{SST k-omega} Backward-Facing-Step Benchmark}
\label{sec:app_A_backstep_ansys_comparison}
The third benchmark evaluates the backward-facing-step geometry introduced in Section~\ref{sec:separated_flow_without_curvature}. This case is included to expose turbulence-model sensitivity in separated internal flow, not to replace the experimental validation evidence from the Driver-Seegmiller/NASA comparisons discussed above. The geometry and boundary conditions are based on the backward-facing-step configuration from Marcib\'{a}l's master thesis and accompanying codebase \cite{Marcibal2024,joove_keps}. They are shown in Figure~\ref{fig:backstep_geometry_bcs}. The flow enters a channel of height $H_\mathrm{in}=2.0$ m with a uniform inlet velocity of $U_\mathrm{in}=25$ m/s. At $x=0$, the lower wall drops by the step height $h=0.25$ m, giving an outlet height of $H_\mathrm{out}=2.25$ m. The benchmark Reynolds number is defined using the step height,
\begin{equation}
	Re_h = \frac{U_\mathrm{in} h}{\nu} = 3.44 \times 10^4 ,
\end{equation}
because backward-facing-step results are conventionally compared using the step height and the reattachment length $x_r/h$. For reference, using the inlet-channel height would give $Re_{H,\mathrm{in}} = 2.75 \times 10^5$, but this is not the primary scaling for the separated region.

\begin{figure}[htpb]
	\centering
	\includegraphics[width=\textwidth]{BackStep-Geometry-BCs}
	\caption{Backward-facing-step benchmark geometry and boundary conditions, based on the configuration used by Marcib\'{a}l and the associated code implementation \cite{Marcibal2024,joove_keps}. The first 5.0 m of the upper and lower boundaries are treated as symmetry/slip segments, while all remaining solid boundaries are no-slip walls. The streamwise dimension is compressed for readability.}
	\label{fig:backstep_geometry_bcs}
\end{figure}

The Fluent setup mirrors the wall-resolved SA assumptions used throughout this thesis. The inlet velocity is prescribed as $\mathbf{u}=(25,0)$ m/s, the outlet pressure is fixed at $p=0$ Pa, and the no-slip walls enforce $\tilde{\nu}=0$. The inlet modified turbulent viscosity is set to $\tilde{\nu}=0.18449$ m$^2$/s, corresponding to the eddy-viscosity level used in the FEniCS backstep configuration. The important comparison quantity is the reattachment length $x_r/h$, obtained from the downstream lower-wall shear stress changing sign from negative to positive. Secondary diagnostics are the size of the recirculation zone, the pressure recovery, and the near-wall skin-friction distribution.

\subsection{Wall-Resolved Ansys Grid-Sensitivity Study}
\label{subsec:app_A_backstep_mesh_ansys}

The grid-sensitivity study uses the same backward-facing-step geometry and named selections as the benchmark definition above. The no-slip wall boundaries are inflated with a fixed first-layer height of $\Delta y_1=1.30\times10^{-4}$ m. This first-layer height is kept constant on all three grids because the Reynolds number and target wall resolution are unchanged. The mesh refinement is therefore applied through the number of inflation layers, inflation growth rate, and streamwise element sizing. Unlike the attached channel and bend verification cases, this separated-flow benchmark is treated as a grid-sensitivity record rather than a strict verification case, while the SA--SST $k-\omega$ comparison below uses the finest converged mesh for each closure.

\begin{table}[htbp]
	\centering
	\scriptsize
	\caption{Wall-resolved backward-facing-step mesh hierarchy for the Ansys Fluent grid-sensitivity study. The first-layer height is fixed to maintain $y^+\approx1$, while the boundary-layer thickness and streamwise resolution are refined.}
	\label{tab:app_A_backstep_meshes_ansys}
	\renewcommand{\arraystretch}{1.2}
	\setlength{\tabcolsep}{5pt}
	\begin{tabular*}{\textwidth}{@{\extracolsep{\fill}}lccc@{}}
		\hline
		\textbf{Parameter} & \textbf{Coarse} & \textbf{Medium} & \textbf{Fine} \\
		\hline
		\multicolumn{4}{@{}l}{\textit{Mesh settings}} \\
		Mesh method & Quad-dom. & Quad-dom. & Quad-dom. \\
		Element order & Linear & Linear & Linear \\
		Inflation option & First layer & First layer & First layer \\
		Inflated boundaries & No-slip walls & No-slip walls & No-slip walls \\
		$\Delta y_1$ [m] & $1.30 \times 10^{-4}$ & $1.30 \times 10^{-4}$ & $1.30 \times 10^{-4}$ \\
			Inflation layers & $28$ & $36$ & $44$ \\
			Growth rate & $1.18$ & $1.15$ & $1.12$ \\
			Near-step wall size [m] & $2.50 \times 10^{-2}$ & $1.25 \times 10^{-2}$ & $6.25 \times 10^{-3}$ \\
			Top-wall size [m] & $5.00 \times 10^{-2}$ & $2.50 \times 10^{-2}$ & $1.25 \times 10^{-2}$ \\
			Inlet/outlet/symmetry size [m] & $1.00 \times 10^{-1}$ & $5.00 \times 10^{-2}$ & $2.50 \times 10^{-2}$ \\
			Nodes & $77{,}725$ & $193{,}276$ & $457{,}175$ \\
			Elements & $86{,}371$ & $192{,}131$ & $454{,}431$ \\
		\hline
		\multicolumn{4}{@{}l}{\textit{Grid metrics: Spalart--Allmaras}} \\
			$x_r/h$ & $5.92$ & $5.44$ & $5.69$ \\
			$\epsilon_{\mathrm{rel}}(x_r/h)$ [\%] & -- & $8.11$ & $4.39$ \\
		Total pressure loss $\Delta p_t$ [Pa] & $46.743$ & $46.795$ & $46.531$ \\
		$\epsilon_{\mathrm{rel}}(\Delta p_t)$ [\%] & -- & $0.11$ & $0.57$ \\
		\multicolumn{4}{@{}l}{\textit{Grid metrics: SST $k-\omega$}} \\
		$x_r/h$ & $6.09$ & $6.13$ & $6.08$ \\
		$\epsilon_{\mathrm{rel}}(x_r/h)$ [\%] & -- & $0.66$& $0.82$ \\
		Total pressure loss $\Delta p_t$ [Pa] & $44.741$ & $45.139$ & $45.142$ \\
		$\epsilon_{\mathrm{rel}}(\Delta p_t)$ [\%] & -- & $0.89$ & $0.01$ \\
		\hline
	\end{tabular*}
	\vspace{2pt}
	\begin{minipage}{0.98\textwidth}
			\scriptsize \textit{Note.} Quad-dom. = quad-dominant. The near-step wall size is applied to \texttt{upstream\_lower\_wall}, \texttt{step\_wall}, and \texttt{downstream\_lower\_wall}. The inflated no-slip walls are \texttt{top\_wall}, \texttt{upstream\_lower\_wall}, \texttt{step\_wall}, and \texttt{downstream\_lower\_wall}; the inlet, outlet, and symmetry boundaries are not inflated. For the $\epsilon_{\mathrm{rel}}$ rows, the Medium column gives the Coarse--Medium change and the Fine column gives the Medium--Fine change. The SA reattachment length is reported as grid-sensitivity evidence rather than proof of strict mesh independence.
	\end{minipage}
\end{table}

The primary grid-sensitivity metric is the normalized reattachment length $x_r/h$, where $x_r$ is obtained from the downstream zero-crossing of the streamwise wall shear stress on the lower wall associated with the end of the primary recirculation region. The SA reattachment length decreases from $x_r/h=5.92$ on the coarse grid to $5.83$ on the medium grid and $5.69$ on the fine grid; the resulting medium--fine change remains above a strict one-percent mesh-independence criterion for this separated-flow metric. The total pressure loss $\Delta p_t$ is retained as a secondary metric and is computed from the mass-flow-weighted inlet-to-outlet total-pressure difference. Velocity profiles are additionally extracted at $x/h=1$, $4$, $6$, and $10$, corresponding to $x=0.25$, $1.00$, $1.50$, and $2.50$ m downstream of the step. These stations cover the initial separated shear layer, the recirculation region, the expected reattachment region, and the downstream recovery region; the SA--SST $k-\omega$ comparison is made using the finest converged mesh available for each turbulence closure.



\FloatBarrier

The contour plots below should therefore be read as a model-form comparison on the same local backward-facing-step geometry, not as a new validation against the NASA Driver-Seegmiller geometry. The streamwise-velocity fields identify the separated low-speed region and the downstream recovery, while the turbulent-viscosity ratio shows how each closure supplies the modeled momentum diffusion that controls this recovery. The velocity contours use the exported range of each Fluent solution, whereas the turbulent-viscosity plots use a common color scale.

\begin{figure}[htpb]
	\centering
	\begin{subfigure}[b]{\textwidth}
		\centering
		\includegraphics[width=0.98\textwidth]{Backstep-X-Velocity-SSTkomega}
		\caption{SST $k-\omega$ contour for $u_x$.}
		\label{fig:backstep_velocity_komega}
	\end{subfigure}
	\vspace{0.25cm}
	\begin{subfigure}[b]{\textwidth}
		\centering
		\includegraphics[width=0.98\textwidth]{Backstep-X-Velocity-SA}
		\caption{SA contour for $u_x$.}
		\label{fig:backstep_velocity_sa}
	\end{subfigure}
	\caption{Streamwise-velocity contours for the backward-facing-step benchmark. Both closures predict a low-speed or reversed-flow region immediately downstream of the step and gradual recovery farther downstream, but the size and diffusion of this region are visibly model dependent.}
	\label{fig:ansys_backstep_velocity_contours}
\end{figure}

Figure~\ref{fig:ansys_backstep_velocity_contours} shows that the two closures give the same qualitative flow topology: separation at the step, a recirculating near-wall region, and redevelopment downstream. The difference is in the local recovery. The SA solution gives a smoother transition from reversed flow to positive streamwise velocity, whereas the SST $k-\omega$ solution retains a more structured low-speed region near the lower wall and a sharper separated shear layer. This is the expected sensitivity for a backward-facing step: reattachment is governed by shear-layer growth and turbulent transport, so small differences in the modeled eddy viscosity can change the local velocity field even when the global recirculation pattern remains plausible.

\begin{figure}[htpb]
	\centering
	\begin{subfigure}[b]{\textwidth}
		\centering
		\includegraphics[width=0.98\textwidth]{Backstep-Turbulent-Viscosity-Ratio-SSTkomega}
		\caption{SST $k-\omega$ contour for $\mu_t/\mu$.}
		\label{fig:backstep_mut_mu_komega}
	\end{subfigure}
	\vspace{0.25cm}
	\begin{subfigure}[b]{\textwidth}
		\centering
		\includegraphics[width=0.98\textwidth]{Backstep-Turbulent-Viscosity-Ratio-SA}
		\caption{SA contour for $\mu_t/\mu$.}
		\label{fig:backstep_mut_mu_sa}
	\end{subfigure}
	\caption{Turbulent-viscosity-ratio contours for the backward-facing-step benchmark, plotted with a common color scale. SA produces a smoother, more moderate eddy-viscosity field, while SST $k-\omega$ gives a more stratified distribution with stronger variation between the near-wall/reverse-flow region and the outer shear-dominated flow.}
	\label{fig:ansys_backstep_mut_mu}
\end{figure}

The turbulent-viscosity contours in Figure~\ref{fig:ansys_backstep_mut_mu} explain the velocity differences. SA is a one-equation model: it transports a modified eddy-viscosity variable and reconstructs $\mu_t$ mainly from local vorticity and wall-distance effects. Its $\mu_t/\mu$ field is therefore relatively smooth in this separated channel. SST $k-\omega$ solves separate equations for turbulent kinetic energy and specific dissipation rate, with near-wall $k-\omega$ behavior and outer-flow blending. In this case it preserves lower eddy viscosity in the immediate wall-bounded/reversed-flow region while carrying a stronger non-uniform eddy-viscosity distribution through the outer part of the channel. The resulting difference is not a numerical detail; it is the model-form sensitivity that makes backward-facing-step flow a useful warning case for SA-based topology optimization.

\begin{figure}[htpb]
	\centering
	\begin{subfigure}[b]{0.49\textwidth}
		\centering
		\includegraphics[width=\textwidth]{Backstep-Velocity-xh1}
		\caption{$x/h=1$}
		\label{fig:backstep_velocity_xh1}
	\end{subfigure}
	\hfill
	\begin{subfigure}[b]{0.49\textwidth}
		\centering
		\includegraphics[width=\textwidth]{Backstep-Velocity-xh4}
		\caption{$x/h=4$}
		\label{fig:backstep_velocity_xh4}
	\end{subfigure}
	\vspace{0.25cm}
	\begin{subfigure}[b]{0.49\textwidth}
		\centering
		\includegraphics[width=\textwidth]{Backstep-Velocity-xh6}
		\caption{$x/h=6$}
		\label{fig:backstep_velocity_xh6}
	\end{subfigure}
	\hfill
	\begin{subfigure}[b]{0.49\textwidth}
		\centering
		\includegraphics[width=\textwidth]{Backstep-Velocity-xh10}
		\caption{$x/h=10$}
		\label{fig:backstep_velocity_xh10}
	\end{subfigure}
	\caption{Normalized streamwise velocity profiles for the backward-facing-step benchmark at four downstream stations. The comparison isolates model-form sensitivity by running SA and SST $k-\omega$ on the same local backward-facing-step geometry and boundary conditions.}
	\label{fig:backstep_velocity_profiles}
\end{figure}

Figures~\ref{fig:ansys_backstep_velocity_contours} and~\ref{fig:ansys_backstep_mut_mu} provide the full-field contour comparison, while Figure~\ref{fig:backstep_velocity_profiles} shows why this benchmark was included even though it is not a direct experimental validation. The purpose is to compare SA and SST $k-\omega$ on the same separated-flow geometry and determine how strongly the modeled mean recirculation and recovery depend on the selected RANS closure. Near the step, at $x/h=1$ and $x/h=4$, both models produce a low-speed or reversed-flow region close to the lower wall, which is consistent with a modeled recirculation bubble generated by the sudden expansion. However, the velocity profiles in this separated region are visibly model dependent, especially near the lower wall and in the developing shear layer. Farther downstream, at $x/h=6$ and especially $x/h=10$, the near-wall streamwise velocity becomes positive and the two profiles move closer together, consistent with modeled reattachment and redevelopment of the boundary layer.

This case is therefore interpreted as a model-sensitivity study rather than as a validation case. Without a direct comparison to the Driver-Seegmiller experimental data for the exact mesh, inlet condition, and domain definition used here, the Fluent results cannot determine which of the two closures is more accurate for this specific setup. The experimental comparison in Figure~\ref{fig:backstep_sa_sst_tmr} already shows the broader trend: in separated backward-facing-step flow, SST $k-\omega$ gives better local velocity-profile agreement than SA, even though both models may predict a reasonable reattachment length. The present Fluent comparison extends that lesson to the geometry used later in the thesis. It shows that the separated region and recovery are sensitive to the turbulence model, so an SA-based optimization result containing separation cannot be interpreted as quantitatively reliable on its own. This limitation is fundamental to linear eddy-viscosity RANS closures: the Boussinesq approximation represents turbulent stresses through a scalar eddy viscosity and cannot fully represent stress anisotropy, non-local turbulence transport, or unsteady coherent structures in separated shear layers. The conclusion is therefore deliberately limited and more conservative than ``SA captures separation'': SA can generate a plausible RANS mean recirculation pattern, but the separated-flow prediction is model dependent and must be checked with an independent closure such as SST $k-\omega$, and ideally with experimental or higher-fidelity validation.

\FloatBarrier
